% 《大学物理简明教程》（第五版）第6章“热力学基础”选择题6.1.1--6.1.13。
% 原书印刷/PDF第153--155页。基于用户上传扫描本独立推导；本章处于待人工二校状态。
\subsection{选择题：平衡、热力学定律与热机}

\sourceproblem{6.1.1}{153}{153}
\problemtext 容器内各处压强相同，或者各处温度相同。这两种单独的条件能否必然说明气体整体处于热力学平衡？
\workings
热力学平衡需要热、力、化学等宏观平衡；单独有$p(\bm r)=\mathrm{const}$并不排除温度梯度与热流，单独有$T(\bm r)=\mathrm{const}$也不保证不存在压强梯度及宏观流动。因此两者均\emph{不必然}是平衡态，原题选\textbf{B}。
\crosscheck 热平衡与力学平衡属于不同条件；不能由一个标量场的均匀性推出整个系统的热力学平衡。

\sourceproblem{6.1.2}{153--154}{153--154}
\problemtext 理想气体内能变化计算式$\Delta U=\nu C_{V,m}\Delta T$在何种过程中适用？其中$\nu$为物质的量、$C_{V,m}$视为常数。
\workings
理想气体的内能是温度的函数，$\mathrm dU=\nu C_{V,m}\mathrm dT$，此关系与经历的路径无关。因此只要气体成分与有效自由度不变，
\[
\boxed{\Delta U=\nu C_{V,m}(T_f-T_i)}
\]
在任意连接两平衡态的过程中均成立，选\textbf{D}（准静态和非准静态均可）。
\crosscheck 这里$\Delta U$为状态函数差而非“热量”；只有等容过程中才另有$Q_V=\Delta U$。

\sourceproblem{6.1.3}{154}{154}
\problemtext 一定量理想气体经历某过程后温度升高，判断系统必定吸热、外界必定对其做正功、内能必定增加，以及“既吸热又对外做正功”是否必然。
\workings
理想气体$U=\nu C_{V,m}T+\mathrm{const}$（热容为正常数），温度升高必有
\[
\boxed{\Delta U>0}.
\]
第一定律$\Delta U=Q-W$并不能确定$Q$或$W$的符号：压缩可以令$W<0$，过程可绝热$Q=0$；加热膨胀则可令$W>0$。只有断言③必然正确，选\textbf{C}。
\crosscheck 对同一内能增量，$Q=\Delta U+W$，只需选取不同的$W$即可构造热流方向不同的过程。

\sourceproblem{6.1.4}{154}{154}
\problemtext 如原图6.1.4，理想气体由同一状态$A(V_1,p_A)$膨胀至同一体积$V_2$，分别经历等压$A\to B$、等温$A\to C$、绝热$A\to D$。问哪一过程吸热最多。
\workings
绝热过程$Q_{AD}=0$。等温过程$\Delta U=0$，故
$Q_{AC}=W_{AC}=\nu RT_A\ln(V_2/V_1)>0$。
等压膨胀过程温度升高，$\Delta U_{AB}>0$且
$W_{AB}=p_A(V_2-V_1)>W_{AC}$，所以
\[
\boxed{Q_{AB}>Q_{AC}>Q_{AD}=0}.
\]
原题选\textbf{A}。
\crosscheck 等压过程的$p(V)$曲线在等温曲线上方；比较同体积区间曲线下面积即可先比较功。

\sourceproblem{6.1.5}{154}{154}
\problemtext 图6.1.5中理想气体从$a$经$b$到$c$，先等容升压再等压膨胀。设气体对外做功为$W$，自外界吸热为$Q$，求$\Delta U,Q,W$符号。
\workings
$a\to b$：$V$不变而$p$升高，故$T$上升、$\Delta U>0$，$W_{ab}=0$。
$b\to c$：$p$不变而$V$增大，故$T$继续上升、$\Delta U>0$，$W_{bc}=p_b(V_c-V_b)>0$。
全程
\[
\boxed{\Delta U>0,\quad W>0,\quad Q=\Delta U+W>0}.
\]
原题选\textbf{B}。
\crosscheck 注意本题$W$特指\emph{气体对外做的功}，不是外界对气体的功。

\sourceproblem{6.1.6}{154}{154}
\problemtext 一台循环热机每循环从$400\,\mathrm K$热源吸收$1800$ J，向$300\,\mathrm K$热源释放$800$ J，对外作功$1000$ J。判断设计能否实现及理由。
\workings
循环的第一定律要求$W=Q_h-Q_c=1000$ J，数字满足能量守恒。
但其热效率
$\eta=W/Q_h=1000/1800=5/9\approx0.5556$；
同温热源间的卡诺效率为
$\eta_{\mathrm C}=1-300/400=0.25$。
\[
\boxed{\eta>\eta_{\mathrm C}},
\]
违反热力学第二定律，所以设计不可实现，选\textbf{D}。
\crosscheck 第一律满足是必要但非充分条件；不能据$W=Q_h-Q_c$就断定热机可行。

\sourceproblem{6.1.7}{154}{154}
\problemtext 理想气体初态温度$T$、体积$V$，依次经历
$(1)$可逆绝热膨胀至$2V$；
$(2)$\emph{等容}升温恢复到$T$；
$(3)$等温压缩回原体积$V$。
判断全循环的净热量、净功和内能改变。
\workings
设气体摩尔数$\nu$、摩尔定容热容$C_{V,m}$、
热容比$\gamma>1$。绝热膨胀末温
$T_2=T\,2^{1-\gamma}<T$，
气体对外作功
\[
W_{12}=\nu C_{V,m}(T-T_2)>0.
\]
第二段\emph{等容}升温，$W_{23}=0$，
$Q_{23}=\nu C_{V,m}(T-T_2)>0$。
第三段在温度$T$下从$2V$等温压缩到$V$，
\[
W_{31}=\nu RT\ln\frac{V}{2V}=-\nu RT\ln2.
\]
总功为
\[
W_{\rm net}=\nu RT\left[
\frac{1-2^{1-\gamma}}{\gamma-1}-\ln2
\right]<0,
\]
这里利用$1-e^{-z}<z$（$z=(\gamma-1)\ln2>0$）。
完整循环状态回到起点，$\Delta U_{\rm cycle}=0$，
由热力学第一定律$Q_{\rm net}=W_{\rm net}<0$。
\[
\boxed{W_{\rm net}<0,\qquad Q_{\rm net}<0,\qquad \Delta U=0}.
\]
气体净向外界放热，原题选\textbf{A}。
\crosscheck 原书第二段确为\textbf{等容变化}而非等压；
故其体积始终为$2V$，没有额外的等压体积功。
该过程的闭合循环$p$--$V$图净面积为负。

\sourceproblem{6.1.8}{154--155}{154--155}
\problemtext 两台可逆卡诺热机的循环面积（$p$--$V$图中包围的面积）相等；第一台工作于温度$T_1$和$T_3$之间，第二台工作于$T_2$与$T_3$之间，$T_1>T_2>T_3$。比较它们的效率、吸热、放热与做功。
\workings
闭合循环对外净功是$p$--$V$闭合曲线面积，
图示面积相等给$W_1=W_2$。
但卡诺效率不同：
\[
\eta_1=1-\frac{T_3}{T_1}>
1-\frac{T_3}{T_2}=\eta_2.
\]
所以$Q_{h,1}=W_1/\eta_1<Q_{h,2}$；两台热机放热也不相等。
唯有
\[
\boxed{Q_{h,1}-Q_{c,1}=Q_{h,2}-Q_{c,2}}
\]
必然成立，选\textbf{D}。
\crosscheck 相等的循环\emph{面积}是相等的净功，而非相等的效率。

\sourceproblem{6.1.9}{155}{155}
\problemtext 关于可逆与准静态过程，判断：①可逆过程必准静态；②准静态过程必可逆；③不可逆过程不能朝相反方向进行；④有摩擦过程必不可逆。
\workings
①真：宏观可逆要求无穷小驱动力且各阶段可视为平衡态；
②假：带有摩擦的缓慢过程也可准静态，却会产生熵；
③假：“不可逆”是不能让系统\emph{连同环境}完全恢复，并非禁止反向变化；
④真：有限摩擦耗散造成不可逆熵产生。
故正确的仅①、④，选\textbf{D}。
\crosscheck 系统可以在外部做功后逆向移动；恢复系统不等于恢复系统与环境的整体。

\sourceproblem{6.1.10}{155}{155}
\problemtext 下面哪个说法严格体现热力学第二定律：从单一热源吸热绝对不可能全部做功、任意可逆过程净吸热等于功、摩擦生热不可逆、热量绝不能从低温传到高温？
\workings
热力学第二定律的开尔文表述禁止\emph{循环工作的装置}只从单一热源吸热且全部转为功，并不排斥非循环的等温膨胀等过程；
任意可逆非循环过程可有$\Delta U\ne0$；
借助外部做功的冰箱可令热从低温流向高温。
摩擦生热会产生熵，无法不留环境改变地完全逆转，故选\textbf{C}。
\crosscheck 断言必须带上“循环”“唯一结果”或“自发”等限定，不能无限推广。

\sourceproblem{6.1.11}{155}{155}
\problemtext 绝热、刚性容器一半为理想气体，另一半为真空，抽去隔板自由膨胀至整个容器，问温度与熵的变化。
\workings
容器绝热且不对外做体积功，$Q=W=0$，故$\Delta U=0$；
理想气体$U=U(T)$，所以
$\boxed{T_f=T_i}$。
自由膨胀不可逆，但可沿同温可逆路径计算熵差：
\[
\boxed{\Delta S=\nu R\ln\frac{V_f}{V_i}
=\nu R\ln2>0}.
\]
温度不变、熵增加，选\textbf{A}。
\crosscheck 不可逆真实路径$Q=0$，\emph{不可}直接沿该路径用$\Delta S=\int\delta Q/T$计算；需任选相同端点间的可逆路径。

\sourceproblem{6.1.12}{155}{155}
\problemtext 对一定量理想气体，哪些组合可能出现：(A)等容加热而内能减少、压强升高；(B)等温压缩同时放热、压强升高；(C)等压压缩时吸热、内能增加；(D)绝热压缩同时内能减少、压强升高。
\workings
等容加热使$T$上升、$U$上升，所以(A)不可能；
等温压缩理想气体$\Delta U=0$，气体对外做负功，$Q=W<0$，而$p\propto1/V$增大，所以(B)可实现。
等压压缩使$T$下降、$\Delta U<0$，故(C)不可能；
绝热压缩时$W<0$、$\Delta U=-W>0$，故(D)不可能。
原题选\textbf{B}。
\crosscheck 所有“加热/放热”与功正负号均以系统\emph{气体}为参考。

\sourceproblem{6.1.13}{155}{155}
\problemtext 图6.1.13中理想气体从$b$压缩到$a$，中间曲线$bca$为绝热路径，
其下方路径为$b1a$，上方路径为$b2a$。
问两条非绝热压缩路径中气体吸放热及做功的符号。
\workings
两路径均由大体积向小体积进行，故
\[
W_{b1a}<0,\qquad W_{b2a}<0.
\]
固定相同端点，$\Delta U=U_a-U_b$相同。
绝热路径$Q_{bca}=0$，且$W_{bca}=-\Delta U$。
下方路径压缩时压力较低，故$|W_{b1a}|<|W_{bca}|$，
$Q_{b1a}=\Delta U+W_{b1a}>0$；
上方路径则$|W_{b2a}|>|W_{bca}|$，故$Q_{b2a}<0$。
\[
\boxed{b1a:\ Q>0,\ W<0;\qquad b2a:\ Q<0,\ W<0}.
\]
原题选\textbf{B}。
\crosscheck 此比较以原图中两条曲线始终分处绝热线两侧为条件；压缩方向会改变路径功符号。
